\chapter{绪论}
\label{chap:introduction}

\section{研究背景与意义}

卷积神经网络（Convolutional Neural Network, CNN）已经成为目标检测、图像分类和语义理解等视觉任务的核心计算模型。YOLO 系列检测器将候选框生成和类别预测整合为单阶段推理流程，在边缘端实时感知任务中具有较好的工程适配性~\cite{redmon2016yolo,redmon2018yolov3}。其中 YOLOv3-tiny 通过减少层数和通道规模降低计算量，适合在嵌入式平台上部署。

边缘端部署面对的约束与数据中心不同。通用 CPU 的吞吐有限，嵌入式 GPU 或专用 NPU 虽然性能较高，但在功耗、可定制性、接口灵活性和可验证性方面并不总是满足定制系统需求。FPGA 与 MPSoC 平台同时包含处理系统（Processing System, PS）和可编程逻辑（Programmable Logic, PL），能够让软件负责图像输入、调度和后处理，让硬件承担规则密集的卷积计算。因此，基于 FPGA 的 CNN 加速器仍然是边缘智能系统中的重要研究方向~\cite{qiu2016going,guo2019angeleye}。

另一方面，真实检测网络的层结构具有明显异构性。以本文采用的 YOLOv3-tiny Conv0--Conv9 链路为例，前端层空间尺寸较大且常伴随池化，后端层空间尺寸较小但通道数显著增加；网络同时包含 3x3 卷积和 1x1 检测头，不同层对 activation、pooling、K pass 数量和输出通道分块的需求也不相同。若硬件仅针对单一卷积核或固定后处理路径设计，则 1x1 层可能被低效映射为稀疏 3x3，池化和激活需要频繁回退到软件，或者每类层都需要独立硬件路径，最终降低资源利用率和系统可扩展性。因此，本文首先需要构建一套层自适应的卷积执行基础，使同一套脉动阵列能够通过运行时描述符适配不同层形状和后处理组合。

然而，仅具备层自适应能力仍不足以保证端到端性能。真实 CNN 网络的加速并不等同于构造一个高峰值 MAC 阵列；每一层还要经过输入通道展开、输出通道分块、空间 tile 分块和多轮 K pass 累加。若数据搬运、部分和反馈或 pass 边界调度处理不当，硬件阵列会在等待权重、等待 IFM replay、等待 partial-PSUM drain 或等待 DMA 完成时产生大量空泡。本文关注的核心问题正是：\textbf{如何在资源受限的 MPSoC 上，构建层自适应的卷积加速基础，并面向真实 YOLOv3-tiny 链式推理降低非计算开销，使优化效果能够被板级实验和性能计数器解释。}

\section{国内外研究现状}

CNN 加速器的研究大致经历了从算子级流水线到空间阵列架构的发展。Kung 提出的 systolic architecture 通过规则 PE 网络在相邻单元之间传递数据，为矩阵乘和卷积计算提供了高复用的硬件组织方式~\cite{kung1982systolic}。Google TPU 使用大规模 systolic array 获得高吞吐推理性能~\cite{jouppi2017tpu}。Eyeriss 系列则从能量效率角度系统讨论了 row-stationary 数据流和层级存储复用~\cite{chen2016eyeriss,chen2019eyerissv2}。Sze 等人的综述对 DNN 加速器的数据流、存储层次和能效权衡进行了系统总结~\cite{sze2017efficient}。

在 FPGA 方向，Qiu 等人较早展示了在嵌入式 FPGA 上部署 CNN 的完整流程~\cite{qiu2016going}；Angel-Eye 进一步提出从模型到 FPGA 映射的设计流程~\cite{guo2019angeleye}；S2A 和 RTL compilation 相关工作关注从高层描述生成 systolic array 或模块化 RTL~\cite{wei2017s2a,ma2018rtlcompiler}。针对 YOLOv3-tiny 的轻量化 FPGA 加速器也已有报道~\cite{ding2022tinyyolo}。这些工作说明 FPGA 对 CNN 推理具有可行性，但许多研究更强调模型映射、阵列生成或单层性能，对真实链式推理中的 DMA 交互、IFM 打包、partial-PSUM 反馈和 pass 级调度开销讨论较少。

本文以可复现的 MPSoC 端到端系统为载体，研究真实 YOLO 链式推理中的数据流和调度瓶颈。相比仅报告理论吞吐或单层仿真的工作，本文更强调硬件 RTL、Vitis runtime、外部 golden、板级 batch-chain 验证、DDR 图像输入和性能计数器之间的一致闭环。

图~\ref{fig:workflow_overview} 给出本文的系统框架与优化路线。论文首先从真实 CNN 链式推理中的数据搬运和调度开销出发，构建 PS/PL 协同的 LASA 卷积执行引擎；随后以 KCS 三维分块模型统一层映射，并通过 BSD、OCRR 和 OPF-P 分别优化 PS/PL 交互、IFM 数据排布和 K-pass 边界空泡；最终通过 golden 验证、DDR image demo 和性能计数器形成板级诊断闭环。

\begin{figure}[htb]
\centering
\includegraphics[width=\textwidth]{\detokenize{img/picture/工作流程图.png}}
\caption{LASA 系统框架与优化路线}
\label{fig:workflow_overview}
\end{figure}

\section{本文主要工作}

围绕上述问题，本文完成的主要工作如下。

\begin{enumerate}
    \item \textbf{提出并实现 LASA 层自适应脉动阵列加速器。} LASA 基于 dual-lane INT8 systolic array，支持 3x3 卷积、原生 1x1 卷积、INT32 partial-PSUM、逐通道重量化、可旁路 activation LUT、可旁路 2x2 maxpool 和 OFM 写回；系统接口采用 AXI-Lite 配置和四路 AXI-Stream/DMA 数据通路，能够与 A53 bare-metal runtime 协同执行 Conv0--Conv9。

    \item \textbf{提出 KCS 三维分块执行模型。} 本文将不同卷积层统一映射为 K-pass、Cout block 和 spatial tile 三个调度维度。K-pass 适配阵列输入 lane 数，Cout block 适配 dual-lane 输出通道并行度，spatial tile 约束片上缓存容量和 DMA 搬运粒度，为后续数据流优化提供统一执行抽象。

    \item \textbf{提出面向非计算开销的三类优化策略。} 面向 PS/PL 交互开销，BSD 通过 batch stream、packet count 和预打包权重减少 DMA 启动与软件服务次数；面向 IFM 布局转换开销，OCRR 通过 raw-HWC tile cache、1x1/3x3 replay 和 backend full-tile 降低软件重排与 tile 边界成本；面向 K-pass 边界气泡，OPF-P 通过流水化 drain、continuous PSUM collector、pass-level prefetch 和 during-compute prefetch 重叠部分和反馈与下一 pass 准备。

    \item \textbf{建立从 RTL 到 KV260 板级运行的验证和性能诊断体系。} 本文使用 Icarus Verilog、Vivado XSIM、外部 RTL golden、链式 semantic golden 和 Vitis board smoke 构成多层验证流程；在硬件中加入 stage、sub-stage、raw-HWC、prefetch、collector、pass timeline 等计数器，用于定位性能瓶颈和解释优化收益。
\end{enumerate}

\section{论文组织结构}

本文共分为六章。第一章介绍研究背景、相关研究现状和本文贡献。第二章给出卷积计算、INT8 量化、systolic array、YOLOv3-tiny 层结构和 MPSoC 平台基础。第三章提出 LASA 层自适应脉动阵列架构，介绍 PS/PL 职责划分、MPSoC 总体结构、LASA 数据通路、层自适应机制和链式推理执行流程。第四章讨论面向链式推理的数据流与调度优化，包括 KCS 三维分块执行模型、BSD 软硬件流式协同、OCRR 片上重排与重放、OPF-P 反馈重叠与预取。第五章介绍运行时集成、验证机制与性能诊断，包括 layer descriptor、STREAM\_CFG 配置、golden 验证路径、性能计数器和最终推荐软硬件组合。第六章给出正确性验证、板级端到端性能、消融实验、资源时序结果和后续工作讨论。
